\chapter{Implementasi dan Evaluasi}

\section{Lingkungan Implementasi}
Bagian ini menjelaskan spesifikasi perangkat keras dan perangkat lunak yang
digunakan untuk implementasi sistem.

\section{Implementasi Antarmuka}
Bagian ini menyajikan tangkapan layar (\emph{screenshot}) hasil implementasi
antarmuka sistem beserta penjelasannya.

\section{Pengujian Sistem}
Pengujian fungsional dilakukan dengan metode \emph{black-box}. Ringkasan hasil
pengujian disajikan pada Tabel~\ref{tab:blackbox}.

\begin{table}[H]
\centering
\caption{Hasil pengujian black-box}
\label{tab:blackbox}
\begin{tabular}{|c|l|l|c|}
\hline
\textbf{Kode} & \textbf{Skenario} & \textbf{Hasil Diharapkan} & \textbf{Status} \\ \hline
\multicolumn{1}{|c|}{(1)} & \multicolumn{1}{c|}{(2)} & \multicolumn{1}{c|}{(3)} & \multicolumn{1}{c|}{(4)} \\ \hline
UJI-01 & Login valid       & Masuk ke beranda          & Sesuai \\
UJI-02 & Unggah data benar & Data tersimpan            & Sesuai \\
UJI-03 & Prediksi          & Hasil prediksi tampil     & Sesuai \\ \hline
\end{tabular}
\end{table}

\section{Evaluasi}
Bagian ini memuat evaluasi sistem, misalnya melalui \emph{User Acceptance Test}
(UAT) atau pengukuran kinerja model, beserta pembahasan hasilnya.
